% GENERATED FILE. Edit canonical lessons, then run npm run generate:books.
% canonical-source-hash: fnv1a64:d8347ec30577ad8f
% canonical-lessons: PA-C285-paratna, PA-C285-varna, PA-C285-mukkna, PA-C285-thakkna, PA-C285-pahu

\chapter{Back Home Again}
\label{ch:pa-a2-back-home-again}
\hlchaptermodality{285}

\begin{chapteropening}
\textbf{By the end of this chapter:} \emph{I can say when I will come back home.}

Complete the last lesson of chapter 285: I can say when I will come back home.
\end{chapteropening}

\section[\texorpdfstring{\pa{ਪਰਤਣਾ} (paratṇā)}{ਪਰਤਣਾ (paratṇā)}]{\pa{ਪਰਤਣਾ} (paratṇā) — to return; to come back}

\label{lesson:PA-C285-paratna}



\begin{quote}
Before the new one: say the Punjabi for last, then the Punjabi for to get to.
\end{quote}

\subsection*{You'll want to know: \pa{ਪਰਤਣਾ}}
\textbf{\pa{ਪਰਤਣਾ}} — \emph{paratṇā} — \textquotedblleft{}to return; to come back\textquotedblright{}.

\begin{quote}
\pa{ਅਸੀਂ} \pa{ਐਤਵਾਰ} \pa{ਨੂੰ} \pa{ਸ਼ਹਿਰ} \pa{ਪਰਤਾਂਗੇ।}
\end{quote}

\begin{quote}
\emph{asīṁ aitvār nū̃ shahir partāṅge.}
\end{quote}

\begin{quote}
We will return to the city on Sunday.
\end{quote}

Past (I …): \pa{ਮੈਂ} \pa{ਘਰ} \pa{ਪਰਤਿਆ}, \emph{maiṁ ghar partiā}.

\begin{cousinweb}
\textbf{\pa{ਪਰਤਣਾ}} (\emph{paratṇā}) is inherited from Sanskrit \emph{parivartatē}, \textquotedblleft{}turns round\textquotedblright{}; CDIAL lists Punjabi \emph{partaṇā} under it only as \textquotedblleft{}to change the side of anything by turning it\textquotedblright{}; the sense \textquotedblleft{}to return\textquotedblright{} is the one it gives for Lahnda \emph{partaṇ} (CDIAL 7872).
\end{cousinweb}

\begin{grammarlens}[title={What you've built}]
One more word for saying what comes next.
\end{grammarlens}

\subsection*{Your turn}
\begin{itemize}
\raggedright
  \item \emph{Say it:} \emph{paratṇā}
  \item \emph{Say it:} \emph{paratṇā}, once more
  \item \emph{Say it:} \emph{asīṁ aitvār nū̃ shahir partāṅge.} — We will return to the city on Sunday.
  \item \emph{Recall:} say the Punjabi for last, then the Punjabi for to get to, then say \emph{paratṇā} again
\end{itemize}

\begin{tcolorbox}[breakable,colback=teal!4,colframe=teal!35!black,title={Before you move on}]
What does \pa{ਪਰਤਣਾ} mean? (\textquotedblleft{}To return; to come back\textquotedblright{}.) Say it once more.
\end{tcolorbox}
\section[\texorpdfstring{\pa{ਵੜਨਾ} (vaṛnā)}{ਵੜਨਾ (vaṛnā)}]{\pa{ਵੜਨਾ} (vaṛnā) — to go in; to enter}

\label{lesson:PA-C285-varna}



\begin{quote}
Before the new one: say the Punjabi for to get to, then the Punjabi for to return.
\end{quote}

\subsection*{You'll want to know: \pa{ਵੜਨਾ}}
\textbf{\pa{ਵੜਨਾ}} — \emph{vaṛnā} — \textquotedblleft{}to go in; to enter\textquotedblright{}.

\begin{quote}
\pa{ਬਿੱਲੀ} \pa{ਰਸੋਈ} \pa{ਵਿੱਚ} \pa{ਵੜਦੀ} \pa{ਹੈ।}
\end{quote}

\begin{quote}
\emph{billī rasoī vich vaṛdī hai.}
\end{quote}

\begin{quote}
The cat goes into the kitchen.
\end{quote}

Past (I …): \pa{ਮੈਂ} \pa{ਕਮਰੇ} \pa{ਵਿੱਚ} \pa{ਵੜਿਆ}, \emph{maiṁ kamre vich vaṛiā}.

\begin{cousinweb}
\textbf{\pa{ਵੜਨਾ}} (\emph{vaṛnā}) is inherited from Sanskrit \emph{upapatati}, \textquotedblleft{}hastens towards\textquotedblright{}; CDIAL lists Punjabi \emph{vaṛnā}, \textquotedblleft{}to enter\textquotedblright{}, under it (CDIAL 2189).
\end{cousinweb}

\begin{grammarlens}[title={What you've built}]
One more word for saying what comes next.
\end{grammarlens}

\subsection*{Your turn}
\begin{itemize}
\raggedright
  \item \emph{Say it:} \emph{vaṛnā}
  \item \emph{Say it:} \emph{vaṛnā}, once more
  \item \emph{Say it:} \emph{billī rasoī vich vaṛdī hai.} — The cat goes into the kitchen.
  \item \emph{Recall:} say the Punjabi for to get to, then the Punjabi for to return, then say \emph{vaṛnā} again
  \item \emph{Write it:} copy \textbf{\pa{ਵੜਨਾ}} once, checking the whole word
\end{itemize}

\begin{tcolorbox}[breakable,colback=teal!4,colframe=teal!35!black,title={Before you move on}]
What does \pa{ਵੜਨਾ} mean? (\textquotedblleft{}To go in; to enter\textquotedblright{}.) Say it once more.
\end{tcolorbox}
\section[\texorpdfstring{\pa{ਮੁੱਕਣਾ} (mukkṇā)}{ਮੁੱਕਣਾ (mukkṇā)}]{\pa{ਮੁੱਕਣਾ} (mukkṇā) — to come to an end; to run out}

\label{lesson:PA-C285-mukkna}



\begin{quote}
Before the new one: say the Punjabi for to return, then the Punjabi for to go in.
\end{quote}

\subsection*{You'll want to know: \pa{ਮੁੱਕਣਾ}}
\textbf{\pa{ਮੁੱਕਣਾ}} — \emph{mukkṇā} — \textquotedblleft{}to come to an end; to run out\textquotedblright{}.

\begin{quote}
\pa{ਛੁੱਟੀ} \pa{ਪਰਸੋਂ} \pa{ਮੁੱਕੇਗੀ।}
\end{quote}

\begin{quote}
\emph{chhuṭṭī parsõ mukkegī.}
\end{quote}

\begin{quote}
The holiday will end the day after tomorrow.
\end{quote}

Past (it …): \pa{ਉਹ} \pa{ਸਫ਼ਰ} \pa{ਮੁੱਕਿਆ}, \emph{oh safar mukkiā}.

\begin{cousinweb}
\textbf{\pa{ਮੁੱਕਣਾ}} (\emph{mukkṇā}) is inherited from a reconstructed \emph{*mukna}, \textquotedblleft{}loosed\textquotedblright{}; CDIAL lists Punjabi \emph{mukkṇā}, \textquotedblleft{}to come to an end, run short\textquotedblright{}, under it (CDIAL 10157).
\end{cousinweb}

\begin{grammarlens}[title={What you've built}]
One more word for saying what comes next.
\end{grammarlens}

\subsection*{Your turn}
\begin{itemize}
\raggedright
  \item \emph{Say it:} \emph{mukkṇā}
  \item \emph{Say it:} \emph{mukkṇā}, once more
  \item \emph{Say it:} \emph{chhuṭṭī parsõ mukkegī.} — The holiday will end the day after tomorrow.
  \item \emph{Recall:} say the Punjabi for to return, then the Punjabi for to go in, then say \emph{mukkṇā} again
  \item \emph{Write it:} copy \textbf{\pa{ਮੁੱਕਣਾ}} once, checking the whole word
\end{itemize}

\begin{tcolorbox}[breakable,colback=teal!4,colframe=teal!35!black,title={Before you move on}]
What does \pa{ਮੁੱਕਣਾ} mean? (\textquotedblleft{}To come to an end; to run out\textquotedblright{}.) Say it once more.
\end{tcolorbox}
\section[\texorpdfstring{\pa{ਥੱਕਣਾ} (thakkṇā)}{ਥੱਕਣਾ (thakkṇā)}]{\pa{ਥੱਕਣਾ} (thakkṇā) — to get tired}

\label{lesson:PA-C285-thakkna}



\begin{quote}
Before the new one: say the Punjabi for to go in, then the Punjabi for to come to an end.
\end{quote}

\subsection*{You'll want to know: \pa{ਥੱਕਣਾ}}
\textbf{\pa{ਥੱਕਣਾ}} — \emph{thakkṇā} — \textquotedblleft{}to get tired\textquotedblright{}.

\begin{quote}
\pa{ਬੱਚੇ} \pa{ਲੰਮੇ} \pa{ਸਫ਼ਰ} \pa{ਵਿੱਚ} \pa{ਥੱਕਦੇ} \pa{ਹਨ।}
\end{quote}

\begin{quote}
\emph{bachche lamme safar vich thakkde han.}
\end{quote}

\begin{quote}
Children get tired on a long journey.
\end{quote}

Past (I …): \pa{ਮੈਂ} \pa{ਬਹੁਤ} \pa{ਥੱਕਿਆ}, \emph{maiṁ bahut thakkiā}.

\begin{cousinweb}
\textbf{\pa{ਥੱਕਣਾ}} (\emph{thakkṇā}) is inherited from a reconstructed \emph{*sthakk}, \textquotedblleft{}stop, halt\textquotedblright{}; CDIAL lists Punjabi \emph{thakkṇā} under it only as \textquotedblleft{}stop, halt\textquotedblright{}; the sense \textquotedblleft{}to become weary\textquotedblright{} is the one it gives for Hindi \emph{thaknā} (CDIAL 13737).
\end{cousinweb}

\begin{grammarlens}[title={What you've built}]
One more word for saying what comes next.
\end{grammarlens}

\subsection*{Your turn}
\begin{itemize}
\raggedright
  \item \emph{Say it:} \emph{thakkṇā}
  \item \emph{Say it:} \emph{thakkṇā}, once more
  \item \emph{Say it:} \emph{bachche lamme safar vich thakkde han.} — Children get tired on a long journey.
  \item \emph{Recall:} say the Punjabi for to go in, then the Punjabi for to come to an end, then say \emph{thakkṇā} again
  \item \emph{Write it:} copy \textbf{\pa{ਥੱਕਣਾ}} once, checking the whole word
\end{itemize}

\begin{tcolorbox}[breakable,colback=teal!4,colframe=teal!35!black,title={Before you move on}]
What does \pa{ਥੱਕਣਾ} mean? (\textquotedblleft{}To get tired\textquotedblright{}.) Say it once more.
\end{tcolorbox}
\section[\texorpdfstring{\pa{ਪਹੁ} (pahu)}{ਪਹੁ (pahu)}]{\pa{ਪਹੁ} (pahu) — dawn; first light}

\label{lesson:PA-C285-pahu}



\begin{quote}
Before the new one: say the Punjabi for to come to an end, then the Punjabi for to get tired.
\end{quote}

\subsection*{You'll want to know: \pa{ਪਹੁ}}
\textbf{\pa{ਪਹੁ}} — \emph{pahu} — \textquotedblleft{}dawn; first light\textquotedblright{}.

\begin{quote}
\pa{ਅਸੀਂ} \pa{ਪਹੁ} \pa{ਤੋਂ} \pa{ਪਹਿਲਾਂ} \pa{ਘਰ} \pa{ਪਹੁੰਚਾਂਗੇ।}
\end{quote}

\begin{quote}
\emph{asīṁ pahu ton pahilāṁ ghar pahunchāṅge.}
\end{quote}

\begin{quote}
We will reach home before dawn.
\end{quote}

A feminine noun.

\begin{cousinweb}
\textbf{\pa{ਪਹੁ}} (\emph{pahu}) is inherited from Sanskrit \emph{prabhā}, \textquotedblleft{}light\textquotedblright{}; CDIAL lists Punjabi \emph{pauh}, \textquotedblleft{}dawn\textquotedblright{}, under it (CDIAL 8705).
\end{cousinweb}

\begin{grammarlens}[title={What you've built}]
One more word for saying what comes next.
\end{grammarlens}

\subsection*{Your turn}
\begin{itemize}
\raggedright
  \item \emph{Say it:} \emph{pahu}
  \item \emph{Say it:} \emph{pahu}, once more
  \item \emph{Say it:} \emph{asīṁ pahu ton pahilāṁ ghar pahunchāṅge.} — We will reach home before dawn.
  \item \emph{Recall:} say the Punjabi for to come to an end, then the Punjabi for to get tired, then say \emph{pahu} again
  \item \emph{Write it:} copy \textbf{\pa{ਪਹੁ}} once, checking the whole word
\end{itemize}

\begin{tcolorbox}[breakable,colback=teal!4,colframe=teal!35!black,title={Before you move on}]
What does \pa{ਪਹੁ} mean? (\textquotedblleft{}Dawn; first light\textquotedblright{}.) Say it once more.
\end{tcolorbox}
